\documentclass[11pt]{article}
\usepackage[margin=1.1in]{geometry}
\usepackage{amsmath,amssymb,amsthm}
\usepackage{hyperref}

\theoremstyle{plain}
\newtheorem{theorem}{Theorem}[section]
\newtheorem{proposition}[theorem]{Proposition}
\newtheorem{lemma}[theorem]{Lemma}
\newtheorem{corollary}[theorem]{Corollary}
\newtheorem{conjecture}[theorem]{Conjecture}
\newtheorem{problem}[theorem]{Problem}
\theoremstyle{definition}
\newtheorem{definition}[theorem]{Definition}
\newtheorem{remark}[theorem]{Remark}

\newcommand{\C}{\mathbb{C}}
\newcommand{\R}{\mathbb{R}}
\newcommand{\Hess}{\operatorname{Hess}}
\newcommand{\adj}{\operatorname{adj}}
\newcommand{\Jac}{\operatorname{Jac}}
\newcommand{\HC}{\mathrm{HC}}
\newcommand{\JC}{\mathrm{JC}}
\newcommand{\Eu}{\mathcal{E}}
\newcommand{\PS}{\mathcal{P}}
\newcommand{\FF}{\mathbb{F}}

\title{Vanishing bordered Hessians: a rigidity theorem, and the exact
localization of the planar Jacobian conjecture inside the four-variable
Hessian conjecture}
\author{HC$_4$ research project\thanks{Companion manuscript to
\cite{companion}. All machine verifications quoted below are exact
(rational/symbolic, sympy 1.14) fail-closed certificate scripts in the
project directory \texttt{certificates/}; see
Appendix~\ref{sec:verification}. Novelty caveat: an adversarial prior-art
audit (recorded in \texttt{prior\_art.md}) demoted several ingredients to
``known, cited'' status --- see Section~\ref{sec:attribution} --- and
located no prior statement of Theorem~\ref{thm:G} or of the localization
Corollary~\ref{cor:headline}; expert referees should nevertheless check
these against the unpublished folklore of the polynomial-automorphism and
affine-differential-geometry communities.}}
\date{August 24, 2026}

\begin{document}
\maketitle

\begin{abstract}
For $e\in\C[x_1,\dots,x_n]$ let
$B(e):=\nabla e^{\top}\adj(\Hess e)\,\nabla e$; up to sign this is the
bordered Hessian determinant, and (over $\R$, at regular points)
$B(e)=0$ says that the level hypersurface of $e$ has vanishing
Gauss--Kronecker curvature. We prove a dimension-free rigidity theorem:
\emph{if $B(e)\equiv0$ then $\det\Hess e\equiv0$}
(Theorem~\ref{thm:G}); the converse is false. Two independent proofs are
given, one analytic (through the local Legendre transform and an Euler
ordinary differential equation on the target) and one purely algebraic
(through Puiseux series at infinity), the latter valid over every field of
characteristic zero. Combining the theorem with the Gordan--Noether
theorem we obtain that every polynomial in at most three variables with
identically vanishing bordered Hessian is, after an affine change of
coordinates, a polynomial in at most two variables
(Corollary~\ref{cor:planar}); this fails from four variables on. We then
apply this to Meng's Hessian conjecture $\HC_4$ --- the statement that
every $f\in\C[x_1,\dots,x_4]$ with $\det\Hess f\in\C^{\times}$ has
polynomial-automorphism gradient --- which, together with the planar
Jacobian conjecture $\JC_2$, is one of the two statements of the
Jacobian/Hessian family left open by the 2026 counterexamples. Call
$v\neq0$ a \emph{pivot} of $f$ if $D_v^2f$ is constant. We prove a
trichotomy (Theorem~\ref{thm:tri}): a four-variable potential with
constant nonzero Hessian determinant that admits a pivot either has
polynomial-automorphism gradient outright (two unconditional cases), or
is, in affine coordinates, exactly a Meng doubling of a planar Keller map
plus an inert planar potential. Consequently \emph{$\HC_4$ restricted to
potentials admitting a pivot is equivalent to $\JC_2$}
(Corollary~\ref{cor:headline}); in particular it holds unconditionally up
to degree five (Corollary~\ref{cor:deg5}, using Moh's theorem).
\textbf{$\HC_4$ itself remains open}: nothing in this paper touches
potentials admitting no pivot, and we state that class as the remaining
problem (Section~\ref{sec:remaining}). Ingredients that turned out to be
known are cited, not claimed: the Euler contraction identity is due to
Nagaoka--Yazawa (and appears in Fox and in Del Pia--Hildebrand--%
Weismantel--Zemmer), the affine equation of the Hessian is Dolgachev's,
and the curvature interpretation of $B$ is due to Reilly and Fox.
\end{abstract}

\section{Introduction}\label{sec:intro}

\subsection{Setting}

Throughout, $\C$ may be replaced by any algebraically closed field of
characteristic zero (and, for Theorem~\ref{thm:G}, by any field of
characteristic zero; see Remark~\ref{rem:fields}). For
$f\in\C[x_1,\dots,x_n]$ write $\Hess f=(\partial^2f/\partial x_i\partial
x_j)$ and let $\adj$ denote the adjugate matrix, so that
$M\adj(M)=\adj(M)M=\det(M)\,I$. A polynomial map $F:\C^n\to\C^n$ with
$\det DF\in\C^{\times}$ is a \emph{Keller map}; the Jacobian conjecture
$\JC_n$ \cite{Keller,BCW} asserts that every Keller map of $\C^n$ is a
polynomial automorphism. Meng's Hessian conjecture $\HC_n$ \cite{Meng2006}
asserts: if $f\in\C[x_1,\dots,x_n]$ has $\det\Hess f\in\C^{\times}$, then
the formal Legendre transform of $f$ is a polynomial --- equivalently,
$\nabla f$ is a polynomial automorphism, and (over $\C$, by the
injectivity-implies-automorphism theorem recalled as
Theorem~\ref{thm:inj} below) equivalently again, $\nabla f$ is injective
\cite{Meng2006,MengYang}.

The July 2026 counterexamples changed the landscape: $\JC_n$ is false for
all $n\ge3$ (Alp\"oge's three-variable map, recorded, verified and
generalized in \cite{MengYang,Gao}), and $\HC_n$ is false for all $n\ge5$
(Meng--Yang \cite{MengYang}). Since $\JC_1$, $\HC_n$ for $n\le3$
\cite{Dillen,deBondt2015} and the bridges $\JC_n\Rightarrow\HC_n$,
$\HC_{2n}\Rightarrow\JC_n$ \cite{Meng2006,MengYang} are settled, exactly
two statements survive: $\JC_2$ and $\HC_4$, linked by
$\HC_4\Rightarrow\JC_2$ (Meng's doubling; the implication is also noted by
de Bondt \cite{deBondt2015}). $\HC_4$ is thus the strongest open statement
of the family. \textbf{This paper does not settle $\HC_4$; it remains
open.} What we prove is an exact description of \emph{where} inside
$\HC_4$ the planar Jacobian conjecture sits.

\subsection{The bordered Hessian and the main results}

\begin{definition}\label{def:B}
For $e\in\C[x_1,\dots,x_n]$ set
\[
B(e):=\nabla e^{\top}\adj(\Hess e)\,\nabla e
\;=\;-\det\begin{pmatrix}\Hess e&\nabla e\\ \nabla e^{\top}&0\end{pmatrix},
\]
the second equality being Cauchy's bordered-determinant expansion
$\det\left(\begin{smallmatrix}M&w\\ w^{\top}&0\end{smallmatrix}\right)
=-w^{\top}\adj(M)\,w$.
\end{definition}

Classically, $B(e)$ computes the Gauss--Kronecker curvature of the level
hypersurfaces of $e$: Reilly \cite[Prop.~4]{Reilly} and Fox
\cite{Fox2017} (whose invariant $U(F)=U^{ij}F_iF_j$ is exactly our $B$,
and who proves $U(F)=\mathcal K\,|dF|^{n+2}$ with $\mathcal K$ the
Euclidean Gauss--Kronecker curvature of the level set) show that, over
$\R$ and at regular points, $B(e)=0$ holds exactly where the level
hypersurface of $e$ is degenerate. We shall say, loosely, that $e$ with
$B(e)\equiv0$ has \emph{developable} level sets; over $\C$ this is only a
name. Our first and central result is a rigidity theorem for this
degenerate case, valid in every dimension.

\begin{theorem}[rigidity; Theorem~\ref{thm:G} below]
Let $e\in\C[x_1,\dots,x_n]$. If $B(e)\equiv0$ then
$\det\Hess e\equiv0$.
\end{theorem}

The converse fails ($e=\tfrac12x_1^2+x_2$ has $\det\Hess e\equiv0$,
$B(e)\equiv1$), so the theorem is genuinely one-directional. Its interest
over $\C$ is that the affine-geometry literature (Fox, Reilly) works under
the standing hypothesis that the quantity
$\langle\nabla e,(\Hess e)^{-1}\nabla e\rangle=B/\det\Hess e$ does
\emph{not} vanish; over $\R$ with definite Hessian the degenerate case is
vacuous. Over $\C$ it has bite, and it is exactly the case that arises in
$\HC_4$. The theorem is also genuinely about \emph{polynomials}: the
algebraic function $x_1-\sqrt{x_1^2-2x_2}$ has $B\equiv0$ and
$\det\Hess\not\equiv0$ (Remark~\ref{rem:conv}).

Combining Theorem~\ref{thm:G} with the Gordan--Noether theorem (through a
homogenization identity of Dolgachev and an Euler-contraction identity of
Nagaoka--Yazawa; both known, see Section~\ref{sec:attribution}) gives:

\begin{corollary}[planarity; Corollary~\ref{cor:planar} below]
Let $n\in\{2,3\}$ and $e\in\C[x_1,\dots,x_n]$ with $B(e)\equiv0$. Then $e$
is affinely $(n-1)$-variable: there is an affine automorphism $\chi$ of
$\C^n$ with $e\circ\chi\in\C[x_1,\dots,x_{n-1}]$.
\end{corollary}

The conclusion is strictly stronger than $\det\Hess e\equiv0$: by de
Bondt--van den Essen's classification of singular Hessians in three
variables \cite{dBvdE} there are zero-Hessian polynomials that are not
affinely $2$-variable, e.g.\ $h=x_1x_2+x_1^2x_3$ --- and consistently
$B(h)=-x_1^4\not\equiv0$. The corollary is sharp in the dimension: it
fails for every $n\ge4$ (Remark~\ref{rem:dim}).

The application to $\HC_4$ goes through the following notion from the
companion manuscript \cite{companion}.

\begin{definition}\label{def:pivot}
Let $f\in\C[x_1,\dots,x_4]$ with $\det\Hess f\in\C^{\times}$. A
\emph{pivot} of $f$ is a vector $v\in\C^4\setminus\{0\}$ such that
$D_v^2f$ is a constant $\gamma\in\C$; the pivot is \emph{quadratic} if
$\gamma\ne0$ and \emph{affine} if $\gamma=0$. In linear coordinates with
$v=e_4$ this says exactly
\[
f=e_0(x')+x_4\,e_1(x')+\tfrac\gamma2x_4^2,\qquad x'=(x_1,x_2,x_3);
\]
for an affine pivot $e_1=D_vf$ is a polynomial constant along $v$, called
the \emph{pivot coefficient}, and its degree $\deg D_vf$ depends only on
$(f,v)$.
\end{definition}

For an affine pivot, the $x_4^2$-coefficient of the bordered expansion of
$\det\Hess f$ is exactly $-B(e_1)$ (Lemma~\ref{lem:bordered}); this is how
vanishing bordered Hessians enter the Hessian conjecture, and it is the
reason Theorem~\ref{thm:G} closes the affine-pivot class. The structural
result is:

\begin{theorem}[pivot trichotomy; Theorem~\ref{thm:tri} below]
Let $f\in\C[x_1,\dots,x_4]$ with $\det\Hess f\in\C^{\times}$ admit a
pivot $v$, and put $\gamma=D_v^2f$. Then exactly one of the following
holds.
\begin{enumerate}
\item[\rm(i)] $\gamma\ne0$; then $\nabla f$ is a polynomial automorphism
(unconditionally, with no hypothesis on the pivot coefficient).
\item[\rm(ii)] $\gamma=0$ and $\deg D_vf=1$; then $\nabla f$ is a
polynomial automorphism (unconditionally).
\item[\rm(iii)] $\gamma=0$ and $\deg D_vf\ge2$; then, after an affine
change of coordinates,
\[
f=a(x_1,x_2)+x_3\,b(x_1,x_2)+x_4\,e(x_1,x_2),
\]
where $(b,e)$ is a planar Keller map; moreover $\nabla f$ is injective if
and only if $(b,e)$ is.
\end{enumerate}
\end{theorem}

Case (iii) is exactly Meng's doubling of the planar Keller map $(b,e)$
\cite{Meng2006,MengYang} plus an ``inert'' planar potential $a$, which
drops out of the Hessian determinant entirely
(Theorem~\ref{thm:doubling}). Case (ii) is \emph{not} in general a
doubling (Remark~\ref{rem:notdoubling}), but it is settled. The headline
consequence:

\begin{corollary}[localization; Corollary~\ref{cor:headline} below]
The following are equivalent:
\begin{enumerate}
\item[\rm(a)] $\JC_2$: every Keller map of $\C^2$ is a polynomial
automorphism;
\item[\rm(b)] every $f\in\C[x_1,\dots,x_4]$ with $\det\Hess f\in
\C^{\times}$ admitting a pivot has $\nabla f$ a polynomial automorphism.
\end{enumerate}
\end{corollary}

Together with Meng's bridge $\HC_4\Rightarrow\JC_2$ this is an exact
localization: the pivot class of $\HC_4$ carries precisely the content of
$\JC_2$, no more and no less, and the only possible source of a
counterexample to $\HC_4$ that is not already a counterexample to $\JC_2$
is the class of potentials with \emph{no} pivot at all
(Section~\ref{sec:remaining}). Since Moh \cite{Moh} proved $\JC_2$ up to
degree $100$, the localization also yields an unconditional statement:
$\HC_4$ holds for every potential of degree at most five that admits a
pivot (Corollary~\ref{cor:deg5}); degree at most four needs no pivot
hypothesis at all, by the companion's degree-four theorem
\cite{companion}.

\subsection{What is claimed and what is not}

Theorem~\ref{thm:G}, Corollary~\ref{cor:planar} as stated for
$B(e)\equiv0$ alone, Theorem~\ref{thm:tri} and
Corollary~\ref{cor:headline} are, to the best of an adversarial
literature search, new; we claim none of them as definitively new pending
expert review. The Euler contraction (Lemma~\ref{lem:identE}) is due to
Nagaoka--Yazawa \cite{NY}; the affine equation of the Hessian
(Lemma~\ref{lem:identH}) is Dolgachev's \cite{Dolgachev}; the curvature
reading of $B$ is Reilly's and Fox's \cite{Reilly,Fox2017};
Theorem~\ref{thm:F} at $n=3$ is essentially a corollary of de Bondt--van
den Essen \cite{dBvdE}; the doubling computation
(Theorem~\ref{thm:doubling}) is elementary given Meng \cite{Meng2006}.
Full attribution is in Section~\ref{sec:attribution}. Every algebraic
identity used below is additionally verified by an exact, fail-closed
machine certificate (Appendix~\ref{sec:verification}); the two proofs of
Theorem~\ref{thm:G} are hand proofs and are \emph{not} machine proofs.

\section{Known inputs}\label{sec:inputs}

We collect the external theorems we use, with sources as verified in the
project's literature files.

\begin{theorem}[Gordan--Noether \cite{GN}; modern proof
\cite{WdB}]\label{thm:GN}
Let $h$ be a homogeneous form in $n\le4$ variables over an algebraically
closed field of characteristic $0$ with $\det\Hess h\equiv0$. Then $h$ is
a cone: its partial derivatives are linearly dependent, i.e.\ there is
$v\ne0$ with $D_vh\equiv0$, and after a linear change of variables $h$
involves at most $n-1$ variables. For $n=2$ this is the classical binary
case: $h=(a_1x_1+a_2x_2)^d$. The statement is false for $n\ge5$ (Perazzo's
cubic $x_1^2x_3+x_1x_2x_4+x_2^2x_5$; see \cite{WdB}).
\end{theorem}

\begin{theorem}[Dillen \cite{Dillen} for $n\le2$; de Bondt
\cite{deBondt2015}, Thm.~2.1, for $n\le3$]\label{thm:hc3}
Let $K$ be a field of characteristic $0$, $n\le3$, and
$f\in K[x_1,\dots,x_n]$ with $\det\Hess f\in K^{\times}$. Then $\nabla f$
is invertible with polynomial inverse. In particular $\nabla f$ is
injective on $K^n$.
\end{theorem}

\begin{theorem}[Bia{\l}ynicki-Birula--Rosenlicht \cite{BBR}; Cynk--Rusek
\cite{CynkRusek}, Thm.~2.2; see also \cite{BCW}, Thm.~2.1]\label{thm:inj}
An injective polynomial map $\C^n\to\C^n$ is bijective and its
set-theoretic inverse is a polynomial map.
\end{theorem}

\begin{theorem}[Moh \cite{Moh}]\label{thm:moh}
$\JC_2$ holds for all planar Keller maps $(f,g)$ over an algebraically
closed field of characteristic $0$ with $\max(\deg f,\deg g)\le100$.
\end{theorem}

For context (not used below): Wang \cite{Wang} proved that every Keller
map of degree $\le2$, in any dimension and any characteristic $\ne2$, is
invertible; Oda's midpoint proof is reproduced in \cite[Thm.~2.4]{BCW}.

The next two lemmas are \emph{known identities}; we state them with their
short proofs only for self-containedness.

\begin{lemma}[Euler contraction; Nagaoka--Yazawa \cite{NY}, Prop.~2.3;
see also \cite{Fox2017}, \cite{DHWZ}, Lemma 5.2, and
\cite{Dolgachev}]\label{lem:identE}
For $e$ homogeneous of degree $d\ge2$ in any number of variables,
\[
B(e)=\nabla e^{\top}\adj(\Hess e)\,\nabla e
=\frac{d}{d-1}\,e\,\det\Hess e .
\]
\end{lemma}

\begin{proof}[Proof (for completeness)]
Write $H=\Hess e$, $g=\nabla e$. Differentiating Euler's relation
$\langle g,x\rangle=de$ gives $Hx=(d-1)g$. Then, using
$H\adj(H)=\det(H)\,I$,
\[
g^{\top}\adj(H)g=(d-1)^{-2}x^{\top}H\adj(H)Hx
=(d-1)^{-2}\det(H)\,x^{\top}Hx
=(d-1)^{-1}\det(H)\,\langle g,x\rangle
=\tfrac{d}{d-1}\,e\det H. \qedhere
\]
\end{proof}

The precise provenance (this is the $s$-linear coefficient of a
proposition of \cite{NY}, machine-verified) is given in
Section~\ref{sec:attribution}.

\begin{lemma}[affine equation of the Hessian; Dolgachev
\cite{Dolgachev}, \S1.1.4, eqs.~(1.20)--(1.21)]\label{lem:identH}
Let $e\in\C[x_1,\dots,x_n]$ be of degree $d\ge2$ and let
$E=x_0^d\,e(x/x_0)$ be its degree-$d$ homogenization. Then
\[
\det\Hess_{n+1}(E)\big|_{x_0=1}
=d(d-1)\,e\,\det\Hess_n(e)-(d-1)^2\,B(e).
\]
\end{lemma}

\begin{proof}[Proof sketch (Dolgachev's)]
At $x_0=1$ Euler's relation for $E$ and for its partials gives
$\partial_0\partial_iE=(d-1)\partial_ie-(\Hess_n(e)\,x)_i$ and
$\partial_0^2E=d(d-1)e-2(d-1)\,x\!\cdot\!\nabla e+x^{\top}\Hess_n(e)\,x$.
Adding $x_i$ times the $i$th row to the $0$th row for $i=1,\dots,n$, then
the mirrored column operation, reduces $\Hess_{n+1}(E)|_{x_0=1}$ to
$\left(\begin{smallmatrix}d(d-1)e&(d-1)\nabla e^{\top}\\
(d-1)\nabla e&\Hess_n(e)\end{smallmatrix}\right)$, and Cauchy's bordered
expansion finishes. (Machine-certified for fully generic $e$ of degrees
$2$ and $3$; see Appendix~\ref{sec:verification}.)
\end{proof}

Finally, Meng's doubling \cite{Meng2006} (see also
\cite[Prop.~2.3]{MengYang}): for a polynomial map $F:\C^n\to\C^n$, the
potential $\varphi(x,y)=\langle y,F(x)\rangle$ on $\C^{2n}$ satisfies
$\det\Hess\varphi=(-1)^n(\Jac F)^2$; this is the mechanism behind
$\HC_{2n}\Rightarrow\JC_n$ and behind Theorem~\ref{thm:doubling} below.

\section{The rigidity theorem}\label{sec:G}

\begin{theorem}[rigidity]\label{thm:G}
Let $n\ge1$ and $e\in\C[x_1,\dots,x_n]$. If $B(e)\equiv0$, then
$\det\Hess e\equiv0$.
\end{theorem}

We give two independent proofs. The first is analytic and exhibits the
mechanism (the inverse gradient would have to be affine in $t^{-1}$
along rays of the target); the second is purely algebraic and works over
any field of characteristic zero. Both begin with the same reduction.

\subsection*{Common set-up}

For $n=1$ the adjugate of the $1\times1$ matrix $(e'')$ is $(1)$, so
$B(e)=(e')^2$; $B\equiv0$ forces $e'\equiv0$, hence $e''\equiv0$. Assume
$n\ge2$, and suppose for contradiction that
\[
B(e)\equiv0\qquad\text{and}\qquad D:=\det\Hess e\not\equiv0 .
\]
Write $g:=\nabla e$, $H:=\Hess e$, and $\varphi:=\nabla e:\C^n\to\C^n$.
Then $e$ is nonconstant and $g\not\equiv0$ (otherwise $H\equiv0$ and
$D\equiv0$).

\begin{proof}[First proof of Theorem~\ref{thm:G} (analytic)]
\emph{Step 1 (base point).} The set $\{x\in\C^n:D(x)\ne0\}$ is a
nonempty Zariski-open subset of $\C^n$, and $g$ does not vanish
identically on it (a polynomial map vanishing on a nonempty Zariski-open
set vanishes identically, which would force $H\equiv0$). Fix $x_0$ with
\[
D(x_0)\ne0
\qquad\text{and}\qquad
p_0:=\varphi(x_0)\ne0 .
\]
These two open conditions are the only genericity used anywhere.

\emph{Step 2 (local inverse and Legendre transform).} Since
$\det D\varphi(x_0)=D(x_0)\ne0$, the holomorphic inverse function theorem
provides a connected open $U\ni p_0$ and a holomorphic $\psi:U\to\C^n$
with $\psi(p_0)=x_0$, $\varphi(\psi(p))=p$ on $U$, and
$D\psi(p)=H(\psi(p))^{-1}$; shrinking $U$ we may assume $D\ne0$ on
$\psi(U)$. Define the local Legendre transform
\begin{equation}\label{eq:legendre}
L(p):=\langle\psi(p),p\rangle-e(\psi(p)),\qquad p\in U .
\end{equation}
The chain rule together with $\varphi\circ\psi=\mathrm{id}_U$ gives
$\nabla L=\psi$ on $U$:
$\partial_jL=\psi_j+\sum_i\partial_j\psi_i\,
\bigl(p_i-\partial_ie(\psi(p))\bigr)=\psi_j$.

\emph{Step 3 (the Euler equation).} Let $\Eu$ denote the Euler operator
on holomorphic functions of $p$:
$(\Eu h)(p):=\frac{d}{d\lambda}h(\lambda p)\big|_{\lambda=1}
=\langle\nabla h(p),p\rangle$. From $\nabla L=\psi$,
\[
\Eu L-L=\langle\psi,p\rangle-\bigl(\langle\psi,p\rangle-e(\psi)\bigr)
=e\circ\psi
\qquad\text{on }U .
\]
Moreover $\nabla(e\circ\psi)(p)=D\psi(p)^{\top}\,\nabla e(\psi(p))
=D\psi(p)^{\top}p$, and $D\psi(p)=H(\psi(p))^{-1}$ is symmetric, so
\[
\Eu(e\circ\psi)(p)=p^{\top}H(\psi(p))^{-1}p
=\Bigl(\frac{g^{\top}\adj(H)\,g}{D}\Bigr)(\psi(p))
=\frac{B(e)}{D}\bigl(\psi(p)\bigr)=0,
\]
because $g(\psi(p))=p$ and $B(e)\equiv0$ as a polynomial. Hence
\begin{equation}\label{eq:eulerode}
(\Eu^2-\Eu)\,L=\Eu(e\circ\psi)=0\qquad\text{on }U .
\end{equation}

\emph{Step 4 (a cone chart around the ray through $p_0$).} Since
$p_0\ne0$, fix an index $m$ with $c:=(p_0)_m\ne0$. For $z\in\C^{n-1}$ let
$\sigma(z)\in\C^n$ be the vector whose $m$th coordinate is $c$ and whose
remaining coordinates, in order, are $z_1,\dots,z_{n-1}$; let $z_0$ be
defined by $\sigma(z_0)=p_0$. The map
\[
\Phi:\C^{\times}\times\C^{n-1}\longrightarrow
\{p\in\C^n:p_m\ne0\},\qquad
\Phi(t,z)=t\,\sigma(z),
\]
is a biholomorphism, with inverse
$p\mapsto\bigl(p_m/c,\;(c/p_m)\,\hat p\bigr)$ where $\hat p$ deletes the
$m$th coordinate. It intertwines scaling with the $t$-coordinate,
$\Phi(\lambda t,z)=\lambda\Phi(t,z)$, so for any holomorphic $h$ defined
near $\Phi(t,z)$,
\[
(\Eu h)(\Phi(t,z))=t\,\partial_t\bigl(h\circ\Phi\bigr)(t,z).
\]
Choose a polydisc $\Delta\ni z_0$ and $\varepsilon\in(0,1)$ so small that
$\Phi(\{|t-1|<\varepsilon\}\times\Delta)\subseteq U$ (possible since
$\Phi(1,z_0)=p_0\in U$ and $\Phi$ is continuous). Put
\[
\Omega:=\Phi(\C^{\times}\times\Delta),
\]
a connected open subset of $\C^n$ (a biholomorphic image of a connected
open set), stable under multiplication by $\C^{\times}$, containing the
whole punctured ray $\C^{\times}p_0=\Phi(\C^{\times}\times\{z_0\})$.

\emph{Step 5 (solving the equation; extension of $L$ over the cone).}
Let $G:=L\circ\Phi$ on $\{|t-1|<\varepsilon\}\times\Delta$. By
\eqref{eq:eulerode} and the intertwining relation,
$\bigl((t\partial_t)^2-t\partial_t\bigr)G=t^2\partial_t^2G=0$; since
$t\ne0$, $\partial_t^2G\equiv0$. For each fixed $z\in\Delta$ the function
$t\mapsto G(t,z)$ is holomorphic on the disc $\{|t-1|<\varepsilon\}$ with
vanishing second derivative, hence affine in $t$:
\[
G(t,z)=A(z)+t\,\widetilde B(z),\qquad
A(z):=G(1,z)-\partial_tG(1,z),\quad
\widetilde B(z):=\partial_tG(1,z),
\]
with $A,\widetilde B$ holomorphic on $\Delta$. Define
\[
\widehat L:\Omega\to\C,\qquad
\widehat L(\Phi(t,z)):=A(z)+t\,\widetilde B(z);
\]
this is holomorphic on $\Omega$ because $\Phi$ is a biholomorphism of
$\C^{\times}\times\Delta$ onto $\Omega$, and $\widehat L=L$ on the open
set $\Phi(\{|t-1|<\varepsilon\}\times\Delta)\ni p_0$. Writing
$\mathcal A(\Phi(t,z)):=A(z)$ and
$\mathcal B(\Phi(t,z)):=t\,\widetilde B(z)$, we have
$\widehat L=\mathcal A+\mathcal B$ on $\Omega$ with, for all $q\in\Omega$
and $\lambda\in\C^{\times}$,
\[
\mathcal A(\lambda q)=\mathcal A(q),\qquad
\mathcal B(\lambda q)=\lambda\,\mathcal B(q)
\]
(homogeneity of degrees $0$ and $1$ on the $\C^{\times}$-stable
$\Omega$; indeed $\lambda\Phi(t,z)=\Phi(\lambda t,z)$).

\emph{Step 6 (the identity $\nabla e\circ\widehat\psi=\mathrm{id}$
extends to the cone).} Put $\widehat\psi:=\nabla\widehat L$, a
holomorphic map $\Omega\to\C^n$; near $p_0$,
$\widehat\psi=\nabla L=\psi$. The map
\[
q\;\longmapsto\;\nabla e\bigl(\widehat\psi(q)\bigr)-q
\]
is holomorphic on $\Omega$ and vanishes on a nonempty open subset of the
\emph{connected} set $\Omega$ (near $p_0$, where $\widehat\psi=\psi$ and
$\varphi\circ\psi=\mathrm{id}$). By the identity theorem for holomorphic
functions of several variables it vanishes on all of $\Omega$:
\begin{equation}\label{eq:section}
\nabla e\bigl(\widehat\psi(q)\bigr)=q\qquad\text{for all }q\in\Omega .
\end{equation}
We emphasize what is (and is not) being claimed: $\widehat\psi$ is a
globally defined holomorphic \emph{section} of $\varphi$ over $\Omega$,
given by the explicit formula of Step~5; no assertion is made that
$\varphi$ is invertible near any point of the ray, that the ray avoids
the critical locus $\{D=0\}$, or that $\widehat\psi$ is the analytic
continuation of $\psi$ as an inverse.

\emph{Step 7 (homogeneity of the section, and the contradiction).}
Differentiating $\mathcal A(\lambda q)=\mathcal A(q)$ with respect to $q$
gives $\lambda\,\nabla\mathcal A(\lambda q)=\nabla\mathcal A(q)$, i.e.\
$\nabla\mathcal A(\lambda q)=\lambda^{-1}\nabla\mathcal A(q)$; similarly
$\nabla\mathcal B(\lambda q)=\nabla\mathcal B(q)$. Hence, with the
\emph{constant} vectors
\[
a:=\nabla\mathcal A(p_0),\qquad b:=\nabla\mathcal B(p_0),
\]
we get, for all $t\in\C^{\times}$,
\[
\widehat\psi(t\,p_0)=t^{-1}a+b .
\]
Substituting $q=t\,p_0$ in \eqref{eq:section}:
\[
\nabla e\bigl(t^{-1}a+b\bigr)=t\,p_0
\qquad\text{for all }t\in\C^{\times}.
\]
Each component of the left-hand side is a polynomial (a component of
$\nabla e$) evaluated at an argument affine in $t^{-1}$, hence is a
polynomial in $t^{-1}$: as a Laurent polynomial in $t$ it contains no
monomial $t^k$ with $k\ge1$. The right-hand side is $t\,p_0$ with
$p_0\ne0$. Two Laurent polynomials agreeing on $\C^{\times}$ are equal,
so comparing the coefficients of $t^{1}$ gives $p_0=0$, contradicting the
choice of $x_0$ in Step~1. Hence $D\equiv0$.
\end{proof}

\begin{remark}[on the continuation step]\label{rem:cont}
The load-bearing step is the passage from the local Legendre transform to
the whole punctured ray. Earlier informal versions of this argument
(including one recorded in the project ledger) phrased it as ``$\psi$
analytically continues along the ray, and the identity
$\varphi\circ\psi=\mathrm{id}$ continues with it'', which as stated is a
non-sequitur: $\psi$ is only a local inverse, the ray $\C^{\times}p_0$
may meet the branch locus of $\varphi$, and nothing guarantees a priori
that the ray stays inside the image of $\varphi$. The proof above is
arranged so that no continuation of $\psi$ \emph{as an inverse} is ever
used: the extension $\widehat L$ is given by an explicit closed formula
(affine in $t$) on an explicit connected cone $\Omega$, and the only
identity transported along the ray, namely
$\nabla e\circ\widehat\psi=\mathrm{id}_\Omega$, is transported by the
identity theorem --- $\widehat\psi$ is merely a section of $\varphi$
afterwards. We consider the argument complete as written, but since it is
the paper's central claim we flag exactly these points --- Steps~5--6 ---
for the referee's scrutiny, and we supply an independent second proof
that uses no function theory at all.
\end{remark}

\begin{proof}[Second proof of Theorem~\ref{thm:G} (algebraic)]
Keep the common set-up ($n\ge2$, $B\equiv0$, $D\not\equiv0$).

\emph{Step 1 ($\C(x)$ is algebraic over $\C(g)$).} The polynomials
$g_1,\dots,g_n$ are algebraically independent over $\C$. Indeed, let
$h\in\C[q_1,\dots,q_n]$ be a nonzero polynomial of minimal degree with
$h(g_1,\dots,g_n)=0$. The chain rule gives
$H\cdot(\nabla h)(g(x))=\nabla\bigl(h(g(x))\bigr)=0$, so
$(\nabla h)(g(x))=0$ on the nonempty Zariski-open set $\{D\ne0\}$, hence
identically; by minimality of $\deg h$ each $\partial_ih$ vanishes, so
$h$ is a nonzero constant, absurd. Consequently
$\C(q_1,\dots,q_n)\to\C(x_1,\dots,x_n)$, $q_i\mapsto g_i$, is an
embedding of fields, and since both sides have transcendence degree $n$
over $\C$, $\C(x)$ is an algebraic extension of (the image of) $\C(q)$.

\emph{Step 2 (Puiseux field at infinity).} Let $\FF:=\C(p_1,\dots,p_n)$
with $p_1,\dots,p_n$ new indeterminates, $\overline\FF$ an algebraic
closure, and
\[
\PS:=\bigcup_{m\ge1}\overline\FF\bigl(\bigl(t^{-1/m}\bigr)\bigr),
\]
the field of Puiseux series at $t=\infty$ with coefficients in
$\overline\FF$. By the Newton--Puiseux theorem, valid over any
algebraically closed field of characteristic $0$
\cite[Ch.~IV, \S3]{Walker}, $\PS$ is algebraically closed. The $n$
elements $tp_1,\dots,tp_n\in\PS$ are algebraically independent over $\C$:
if $Q\in\C[q]$ satisfies $Q(tp)=0$, then writing $Q=\sum_kQ_k$ in
homogeneous parts, $Q(tp)=\sum_kt^kQ_k(p)$, so every $Q_k(p)=0$ in $\FF$
and $Q=0$. Hence $q_i\mapsto tp_i$ defines an embedding $\C(q)\to\PS$.
Since $\C(x)$ is algebraic over $\C(q)$ (Step~1) and $\PS$ is
algebraically closed, this embedding extends to an embedding of fields
\cite[Ch.~V, \S2]{Lang}
\[
\iota:\C(x)\longrightarrow\PS,\qquad
\chi:=(\iota(x_1),\dots,\iota(x_n))\in\PS^n .
\]
Two consequences: (E1) $g(\chi)=t\,p$ as vectors in $\PS^n$ (this is
$\iota(g_i)=\iota(q_i)=tp_i$); (E2) $\iota(h)\ne0$ for every nonzero
$h\in\C(x)$, since $\iota$ is a field embedding --- in particular
$D(\chi)\ne0$, so $H(\chi)$ is invertible over $\PS$, and
$B(\chi)=0$ since $B$ is the zero polynomial.

\emph{Step 3 (derivations).} The formal derivative $\partial_t$ acts
$\overline\FF$-linearly on each $\overline\FF((t^{-1/m}))$ (with
$\partial_tt^{k/m}=\tfrac km t^{k/m-1}$), compatibly in $m$, hence on
$\PS$; its kernel is exactly $\overline\FF$ (a series with all
coefficients of $t^{k/m}$, $k\ne0$, killed is a constant). Each
$\partial/\partial p_i$ is a derivation of $\FF$; since $\overline\FF$ is
a separable algebraic extension of $\FF$ (characteristic $0$), it extends
uniquely to a derivation $D_i$ of $\overline\FF$
\cite[Ch.~VIII, \S5]{Lang}, and then coefficientwise to $\PS$ with
$D_i(t^{k/m})=0$. All these derivations commute with each other and
with $\partial_t$, and they satisfy the chain rule for polynomial
expressions: for $P\in\C[x]$ and any derivation $\partial$ of $\PS$,
$\partial\bigl(P(\chi)\bigr)=\sum_j(\partial_jP)(\chi)\,\partial\chi_j$.

\emph{Step 4 (the formal Legendre value is affine in $t$).} Apply
$\partial_t$ to (E1): $H(\chi)\,\partial_t\chi=p$, so
$\partial_t\chi=H(\chi)^{-1}p$. Put
\[
\Lambda:=\langle\chi,tp\rangle-e(\chi)\in\PS .
\]
Then, using (E1),
\[
\partial_t\Lambda=\langle\partial_t\chi,tp\rangle+\langle\chi,p\rangle
-\langle g(\chi),\partial_t\chi\rangle=\langle\chi,p\rangle ,
\]
and, since $p$ is $t$-free and $p=t^{-1}g(\chi)$,
\[
\partial_t^2\Lambda=\langle\partial_t\chi,p\rangle
=p^{\top}H(\chi)^{-1}p
=t^{-2}\,\frac{g^{\top}\adj(H)\,g}{D}(\chi)
=t^{-2}\,\frac{B(\chi)}{D(\chi)}=0 .
\]
As $\ker\partial_t=\overline\FF$, there are $\alpha,\beta\in\overline\FF$
with $\Lambda=\alpha+\beta t$.

\emph{Step 5 (the inverse gradient is affine in $t^{-1}$).} Apply $D_i$
to the definition of $\Lambda$, using $D_i(tp_j)=t\,\delta_{ij}$ and (E1):
\[
D_i\Lambda=\langle D_i\chi,tp\rangle+t\chi_i-\langle g(\chi),D_i\chi\rangle
=t\,\chi_i .
\]
On the other hand $D_i\Lambda=D_i\alpha+t\,D_i\beta$ with
$D_i\alpha,D_i\beta\in\overline\FF$. Hence, with the $t$-free vectors
$a:=(D_1\alpha,\dots,D_n\alpha)$, $b:=(D_1\beta,\dots,D_n\beta)
\in\overline\FF^{\,n}$,
\[
\chi=t^{-1}a+b .
\]

\emph{Step 6 (contradiction).} Substituting into (E1):
$g(t^{-1}a+b)=t\,p$ in $\PS^n$. The left-hand side is a polynomial map
evaluated at a vector affine in $t^{-1}$ with $t$-free coefficients, so
each component lies in $\overline\FF[t^{-1}]$ and involves only
exponents $\le0$; the right-hand side has the term $t^{1}p_i$ with
$p_i\ne0$. Since Puiseux expansions are unique, this is impossible. Hence
$D\equiv0$.
\end{proof}

\begin{remark}[a third, elementary variant]\label{rem:flow}
The two proofs are the same skeleton (the Legendre value is affine in the
ray parameter, so the inverse gradient is affine in its reciprocal, which
a polynomial gradient cannot match). A third variant, found in an
independent audit of the project, stays entirely on the source side: on
$\{D\ne0\}$ put $V:=H^{-1}g$ and $M:=\sum_k(\partial_kH)V_k$. Two lines of
matrix calculus give $\mathrm{Jac}(V)=I-H^{-1}M$ and
$\nabla(g^{\top}V)=2g-MV$ (both certified fully generically, see
Appendix~\ref{sec:verification}); since $g^{\top}V=B/D\equiv0$ one gets
$MV=2g$, hence $\mathrm{Jac}(V)\,V=-V$. The local holomorphic flow
$x'(s)=V(x(s))$, $x(0)=x_0$, therefore satisfies $x''=-x'$, i.e.\
$x(s)=a_0e^{-s}+b_0$, while $\frac{d}{ds}\nabla e(x(s))=HV=g$ gives
$\nabla e(x(s))=e^{s}p_0$. Thus the polynomial vector
$\tau\,\nabla e(a_0\tau+b_0)-p_0$ vanishes at the infinitely many values
$\tau=e^{-s}$ but not at $\tau=0$ --- the same contradiction. (A formal
power-series version avoids analysis.)
\end{remark}

\begin{remark}[hypotheses]\label{rem:hyp}
No hypothesis on $e$ is needed beyond $B(e)\equiv0$: constancy of $e$ is
not excluded (the statement is then trivial), and the two genericity
conditions of Step~1 ($D(x_0)\ne0$, $\nabla e(x_0)\ne0$) are open
conditions automatically satisfiable under the contradiction hypothesis
$D\not\equiv0$. The first proof uses the holomorphic inverse function
theorem and the identity theorem and is carried out over $\C$; the second
uses only field theory.
\end{remark}

\begin{remark}[other fields]\label{rem:fields}
The second proof is valid verbatim over any field $k$ of characteristic
$0$ (only $\overline\FF$ needs to be algebraically closed, and it is).
Independently, the statement descends from $\C$ to any subfield: the
coefficients of a given $e$ generate a finitely generated extension of
$\mathbb{Q}$, which embeds into $\C$, and $B(e)$, $\det\Hess e$ are
computed coefficient-functorially.
\end{remark}

\begin{remark}[converse; sharpness; low dimensions; machine
checks]\label{rem:conv}
(i) The converse is false: $e=\tfrac12x_1^2+x_2$ has $\det\Hess
e\equiv0$ and $B(e)\equiv1$.
(ii) Polynomiality is used exactly once (the last step of either proof),
and cannot be dropped: the algebraic function
$F=x_1-\sqrt{x_1^2-2x_2}$ satisfies $B(F)\equiv0$ and
$\det\Hess F=-(x_1^2-2x_2)^{-2}\not\equiv0$, and its local inverse
gradient has exactly the form $t^{-1}a+b$ along rays that the proofs
predict. Likewise every rational function homogeneous of degree $0$ has
$B\equiv0$ by the Euler contraction with $d=0$. Theorem~\ref{thm:G} is
therefore a statement about polynomials, not about $C^\infty$ or
algebraic functions --- which may explain its absence from the
affine-geometry literature.
(iii) For $n=2$ there is a completely elementary independent proof: with
$J=\left(\begin{smallmatrix}0&1\\-1&0\end{smallmatrix}\right)$ one has
$\adj H=J^{\top}HJ$, so $B=T^{\top}HT$ with $T=Jg$ tangent to the level
curves; $B\equiv0$ says every level curve is a union of lines; a generic
fibre avoids the finitely many critical values of $e$, hence is smooth,
so its line components are disjoint, i.e.\ parallel; comparing two
generic fibres forces $e\in\C[\ell]$ for a single linear form $\ell$,
whence $\det\Hess e\equiv0$. (An early machine scan accompanying this
argument dropped the linear terms of $e$ --- on which $B$ depends --- and
was not exhaustive as first claimed; the full-coefficient-space decision
in degrees $\le4$ is \texttt{pzG\_n2\_decision.py}.)
(iv) The certificates \texttt{theorem\_G.py}, \texttt{advG\_flow\_proof.py}
and \texttt{pzG\_audit\_identities.py} verify the structural identities
of the proofs fully symbolically and test the conclusion decisively on
complete low-degree ans\"atze. These are consistency certificates for hand
proofs, not machine proofs of the theorem.
\end{remark}

\section{Planarity in at most three variables}\label{sec:planar}

\begin{theorem}\label{thm:F}
Let $n\in\{2,3\}$ and let $e\in\C[x_1,\dots,x_n]$ satisfy
\[
B(e)\equiv0\qquad\text{and}\qquad\det\Hess_n(e)\equiv0 .
\]
Then $e$ is affinely $(n-1)$-variable: there is an affine automorphism
$\chi$ of $\C^n$ with $e\circ\chi\in\C[x_1,\dots,x_{n-1}]$. Equivalently,
there is $v\in\C^n\setminus\{0\}$ with $D_ve\equiv0$.
\end{theorem}

\begin{proof}
If $d:=\deg e\le1$ the claim is trivial (any affine $e$ is affinely
$(n-1)$-variable for $n\ge2$). Assume $d\ge2$ and let $E$ be the
degree-$d$ homogenization of $e$ in the variables $(x_0,x_1,\dots,x_n)$.
By Lemma~\ref{lem:identH} and the two hypotheses,
$\det\Hess_{n+1}(E)|_{x_0=1}\equiv0$; since $\det\Hess_{n+1}(E)$ is
homogeneous and a homogeneous polynomial vanishing on the affine chart
$\{x_0=1\}$ vanishes identically, $\det\Hess_{n+1}(E)\equiv0$. As
$n+1\le4$, Theorem~\ref{thm:GN} applies: there is
$v=(v_0,v')\ne0$, $v'\in\C^n$, with $D_vE\equiv0$.

\emph{Case $v_0=0$.} Then $v'\ne0$ and restricting $D_{v'}E\equiv0$ to
$x_0=1$ gives $D_{v'}e\equiv0$: $e$ is constant in the direction $v'$,
hence linearly $(n-1)$-variable.

\emph{Case $v_0\ne0$.} Normalize $v_0=1$. Euler's relation for $E$ gives
$\partial_0E|_{x_0=1}=d\,e-x\!\cdot\!\nabla e$, so $D_vE\equiv0$ restricted
to $x_0=1$ reads $d\,e=(x-v')\!\cdot\!\nabla e$. Setting
$\tilde e(y):=e(y+v')$ this becomes $y\!\cdot\!\nabla\tilde e=d\,\tilde e$,
and comparing homogeneous parts ($y\!\cdot\!\nabla\tilde e_k=k\tilde e_k$)
shows $\tilde e$ is homogeneous of degree $d\ge2$. $B$ is translation
invariant ($\Hess$ and $\nabla$ commute with translations), so
$B(\tilde e)\equiv0$; Lemma~\ref{lem:identE} gives
$0=B(\tilde e)=\tfrac{d}{d-1}\tilde e\,\det\Hess\tilde e$, and since
$\tilde e\ne0$ in the domain $\C[y]$, $\det\Hess_n\tilde e\equiv0$. Now
Theorem~\ref{thm:GN} in $n\le3$ variables makes $\tilde e$ a cone:
$D_{v''}\tilde e\equiv0$ for some $v''\ne0$, i.e.\ $D_{v''}e\equiv0$.
Composing the translation by $v'$ with a linear change carrying $v''$ to
$e_n$ gives the affine $\chi$.
\end{proof}

\begin{remark}[attribution for $n=3$]\label{rem:Fknown}
For $n=3$, Theorem~\ref{thm:F} is essentially known: it is a short
corollary of the classification of three-variable polynomials with
identically singular Hessian by de Bondt--van den Essen
\cite[Thm.~3.3]{dBvdE} (restated in \cite{deBondt2015}): after a linear
change such a polynomial is either in $\C[x_1,x_2]+\C x_3$ or of the form
$g_1(x_1)+g_2(x_1)x_2+g_3(x_1)x_3$, and in the second family one computes
$B=-(g_2g_3'-g_3g_2')^2$, so $B\equiv0$ forces $g_2,g_3$ to be linearly
dependent, whence $D_ve\equiv0$ for a suitable $v\in\{0\}\times\C^2$
(certificates \texttt{tgc\_debondt\_chain.py},
\texttt{gcv\_family\_wronskian.py}). The self-contained homogenization
proof above is retained because it treats $n=2$ and $n=3$ uniformly and
keeps the paper independent of the classification.
\end{remark}

\begin{corollary}[planarity]\label{cor:planar}
Let $n\in\{2,3\}$ and $e\in\C[x_1,\dots,x_n]$ with $B(e)\equiv0$. Then
$e$ is affinely $(n-1)$-variable. In particular, for $n=2$: $e=\phi(\ell)$
for some $\phi\in\C[s]$ and some affine-linear $\ell$.
\end{corollary}

\begin{proof}
Theorem~\ref{thm:G} gives $\det\Hess_ne\equiv0$; Theorem~\ref{thm:F}
applies.
\end{proof}

\begin{remark}[the dimension barrier is sharp]\label{rem:dim}
The proof of Theorem~\ref{thm:F} consumes the Gordan--Noether theorem in
$n+1$ variables, available exactly for $n+1\le4$, and the corollary
genuinely fails from $n=4$ on: $e=x_2+x_1x_3+x_1^2x_4$ has $B(e)\equiv0$
(its Hessian has a zero row and rank $2$, so its adjugate vanishes) and
$\det\Hess e\equiv0$, yet its four partial derivatives
$x_3+2x_1x_4,\,1,\,x_1,\,x_1^2$ are linearly independent, so $e$ is not
affinely $3$-variable. Its homogenization $x_0^2x_2+x_0x_1x_3+x_1^2x_4$ is
exactly Perazzo's quinary cubic, the canonical non-cone with vanishing
Hessian: the failure of the corollary in four variables is the failure
of Gordan--Noether in five, dehomogenized. More generally
$B\bigl(a(x_1,x_2,x_3)+\lambda x_4\bigr)=\lambda^2\det\Hess_3(a)$ (a
jet-level identity), so every three-variable zero-Hessian polynomial that
is not affinely $2$-variable yields a four-variable counterexample.
Theorem~\ref{thm:G} itself, by contrast, is dimension-free.
\end{remark}

\begin{remark}[strictness]\label{rem:strict}
$B(e)\equiv0$ is strictly stronger than $\det\Hess e\equiv0$ for $n=3$:
the de Bondt--van den Essen polynomial $h=x_1x_2+x_1^2x_3$ has
$\det\Hess h\equiv0$ but is not affinely $2$-variable, and indeed
$B(h)=-x_1^4\not\equiv0$, as Corollary~\ref{cor:planar} requires.
\end{remark}

\section{Pivots in four variables and the localization of
$\JC_2$}\label{sec:pivot}

Throughout this section $f\in\C[x_1,\dots,x_4]$,
$\det\Hess f=c\in\C^{\times}$, and pivots are as in
Definition~\ref{def:pivot}. An affine change of coordinates
$x\mapsto Ax+x^{\ast}$ multiplies $\det\Hess f$ by $(\det A)^2$, carries
pivots to pivots ($D_{Av}^2f=D_v^2(f\circ\chi)\circ\chi^{-1}$ up to the
obvious identification), preserves $\deg D_vf$, and preserves the
polynomial-automorphism property of the gradient
($\nabla(f\circ\chi)=A^{\top}(\nabla f)\circ\chi$); so we may and do work
in pivot-adapted coordinates.

\subsection{The bordered form of the determinant and the developability
constraint}

\begin{lemma}\label{lem:bordered}
Let $f=e_0(x')+x_4\,e_1(x')$ with $x'=(x_1,x_2,x_3)$, and write
$P:=\Hess_3e_0$, $K:=\Hess_3e_1$, $g:=\nabla'e_1$. Then
\[
\Hess f=\begin{pmatrix}P+x_4K&g\\ g^{\top}&0\end{pmatrix},
\qquad
\det\Hess f=-\,g^{\top}\adj\bigl(P+x_4K\bigr)\,g ,
\]
and, as a polynomial in $x_4$ (of degree $\le2$, since the adjugate of a
$3\times3$ pencil is quadratic in the pencil parameter),
\[
[x_4^2]\;\det\Hess f=-\,g^{\top}\adj(K)\,g=-B(e_1).
\]
Consequently, if $\det\Hess f\in\C^{\times}$ then $B(e_1)\equiv0$, and
$e_1$ is nonconstant (a constant $e_1$ zeroes the last row of $\Hess f$).
\end{lemma}

\begin{proof}
The block form is immediate from $\partial_4^2f=0$,
$\partial_i\partial_4f=\partial_ie_1$; the determinant formula is
Cauchy's bordered expansion (Definition~\ref{def:B}). Writing
$\adj(P+tK)=\adj P+tL(P,K)+t^2\adj K$ (each entry is a $2\times2$ minor
of $P+tK$, hence quadratic in $t$ with $t^2$-coefficient the
corresponding minor of $K$) and extracting the $t^2$-coefficient gives
the claim, since $g=\nabla e_1$ and $K=\Hess e_1$ make
$g^{\top}\adj(K)g=B(e_1)$.
\end{proof}

Geometrically, by Reilly--Fox the constraint $B(e_1)\equiv0$ says that
all level surfaces of the pivot coefficient are developable
\cite{Reilly,Fox2017}; e.g.\ a nondegenerate quadric pivot coefficient
such as $e_1=x_1x_2+x_3$ (with $B(e_1)\equiv1$) is impossible.

\subsection{The two unconditional pivot cases}

\begin{theorem}[quadratic pivot; pivot theorem of \cite{companion},
proof included]\label{thm:quad}
Let $f=e_0(x')+x_4e_1(x')+\tfrac\gamma2x_4^2$ with $\gamma\ne0$ and
$\det\Hess f=c\in\C^{\times}$. Then $\nabla f$ is a polynomial
automorphism --- with no hypothesis on $e_1$.
\end{theorem}

\begin{proof}
Set $\tilde e_0:=e_0-e_1^2/(2\gamma)$, $K:=\Hess_3e_1$,
$g:=\nabla'e_1$, $u:=x_4+e_1(x')/\gamma$. Since
$\Hess_3(e_1^2/2)=e_1K+gg^{\top}$,
\[
\Hess_3\tilde e_0+uK
=\Hess_3e_0-\frac{e_1K+gg^{\top}}{\gamma}+x_4K+\frac{e_1}{\gamma}K
=\bigl(\Hess_3e_0+x_4K\bigr)-\frac{gg^{\top}}{\gamma},
\]
so the Schur complement of the $(4,4)$ entry $\gamma$ in
$\Hess f=\left(\begin{smallmatrix}\Hess_3e_0+x_4K&g\\
g^{\top}&\gamma\end{smallmatrix}\right)$ gives
\[
\det\Hess f=\gamma\,\det\nolimits_3\!\bigl(\Hess_3\tilde e_0(x')
+u\,K(x')\bigr)\equiv c .
\]
Now fix $x'\in\C^3$: $K(x')$ is then a constant matrix, and as $x_4$
sweeps $\C$ so does $u$. Hence the \emph{pencil identity}
\[
\det\nolimits_3\!\bigl(\Hess_3\tilde e_0(x')+s\,K(x')\bigr)=c/\gamma
\qquad\text{for all }s\in\C\text{ and all }x'\in\C^3 .
\]
In particular, for each \emph{scalar} $\lambda\in\C$ the three-variable
polynomial $h_\lambda:=\tilde e_0+\lambda e_1$ has
$\Hess_3h_\lambda=\Hess_3\tilde e_0+\lambda K$ and so
$\det\Hess_3h_\lambda\equiv c/\gamma\in\C^{\times}$; thus
$\nabla'h_\lambda$ is injective by Theorem~\ref{thm:hc3}. (Constancy of
$K$ is never needed: the sweep is per-point in $x'$.)

Collision transfer: suppose $\nabla f(p)=\nabla f(q)$. The fourth
component reads $e_1(p')+\gamma p_4=e_1(q')+\gamma q_4$, so
\[
\lambda:=p_4+e_1(p')/\gamma=q_4+e_1(q')/\gamma .
\]
The first three components read
$\nabla'e_0(p')+p_4g(p')=\nabla'e_0(q')+q_4g(q')$; since
$\nabla'h_\lambda=\nabla'e_0-\tfrac{e_1}{\gamma}g+\lambda g$ and
$\lambda-e_1(p')/\gamma=p_4$ (resp.\ at $q'$), this says precisely
$\nabla'h_\lambda(p')=\nabla'h_\lambda(q')$. Injectivity of
$\nabla'h_\lambda$ gives $p'=q'$, and then $p_4=q_4$ from $\lambda$.
Thus $\nabla f$ is injective, hence a polynomial automorphism by
Theorem~\ref{thm:inj}.
\end{proof}

\begin{theorem}[affine pivot coefficient; pivot theorem of
\cite{companion}, case $\deg e_1\le1$, proof included]\label{thm:aff}
Let $f=e_0(x')+x_4e_1(x')$ with $\det\Hess f=c\in\C^{\times}$ and $e_1$
affine (necessarily nonconstant). Then $\nabla f$ is a polynomial
automorphism --- with no hypothesis on $\deg e_0$.
\end{theorem}

\begin{proof}
Write $e_1=\ell(x')+\delta$ with $\ell\ne0$ linear. Choose $S\in
GL_3(\C)$ with $\ell(Sy)=y_3$ and replace $f$ by
$f\circ(S\oplus\mathrm{id}_{x_4})-\delta x_4$: this changes $\nabla f$ by
an invertible linear substitution and a constant shift, preserving both
$\det\Hess f\in\C^{\times}$ and injectivity, so WLOG $e_1=x_3$. By
Lemma~\ref{lem:bordered} with $g=(0,0,1)^{\top}$, $K=0$:
\[
\det\Hess f=-\adj(P)_{33}=-\det\nolimits_2\Hess_{(x_1,x_2)}e_0\equiv c ,
\]
identically in $(x_1,x_2,x_3)$. Hence for every $a\in\C$ the planar
polynomial $e_0(\cdot,\cdot,a)$ has constant Hessian determinant
$-c\in\C^{\times}$, and its gradient is injective by
Theorem~\ref{thm:hc3} ($n=2$, Dillen). Now
$\nabla f=(\partial_1e_0,\;\partial_2e_0,\;\partial_3e_0+x_4,\;x_3)$: a
collision forces equal $x_3$ (fourth component), then equal $(x_1,x_2)$
(first two components, Dillen on the common fibre $x_3=a$), then equal
$x_4$ (third component). So $\nabla f$ is injective, hence an
automorphism by Theorem~\ref{thm:inj}.
\end{proof}

\subsection{The doubling class}

\begin{theorem}[doubling structure; elementary, after Meng
\cite{Meng2006}]\label{thm:doubling}
For any $a,b,e\in\C[x_1,x_2]$, the potential
\[
f=a(x_1,x_2)+x_3\,b(x_1,x_2)+x_4\,e(x_1,x_2)
\]
satisfies, independently of $a$,
\[
\Hess f=\begin{pmatrix}A&M\\ M^{\top}&0\end{pmatrix},\quad
M=\begin{pmatrix}b_{x_1}&e_{x_1}\\ b_{x_2}&e_{x_2}\end{pmatrix},
\qquad
\det\Hess f=\bigl(\Jac(b,e)\bigr)^2 .
\]
In particular $\det\Hess f\in\C^{\times}$ if and only if $(b,e)$ is a
planar Keller map. Thus $f$ is exactly the Meng doubling
$\langle(x_3,x_4),(b,e)\rangle$ plus an inert planar potential $a$.
\end{theorem}

\begin{proof}
The block form is direct differentiation ($f$ is affine in
$(x_3,x_4)$). Moving the last two rows above the first two is an even
permutation of the rows and produces the block-triangular matrix
$\left(\begin{smallmatrix}M^{\top}&0\\ A&M\end{smallmatrix}\right)$,
whence $\det\Hess f=\det(M^{\top})\det(M)=(\Jac(b,e))^2$. Finally,
$(\Jac(b,e))^2\in\C^{\times}$ forces the polynomial $\Jac(b,e)$ to be a
nonzero constant.
\end{proof}

\begin{lemma}[injectivity transfer]\label{lem:inj}
In the situation of Theorem~\ref{thm:doubling}, assume $(b,e)$ is a
Keller map. Then $\nabla f$ is injective if and only if
$(b,e):\C^2\to\C^2$ is injective.
\end{lemma}

\begin{proof}
Write $x'=(x_1,x_2)$ and $M(x')$ as above, invertible for every $x'$
since $\det M=\Jac(b,e)\in\C^{\times}$. The components of $\nabla f$ are
\[
\nabla f=\bigl(a_{x_1}+x_3b_{x_1}+x_4e_{x_1},\;
a_{x_2}+x_3b_{x_2}+x_4e_{x_2},\;b,\;e\bigr).
\]
($\Leftarrow$) If $\nabla f(p)=\nabla f(q)$, components $3,4$ give
$(b,e)(p')=(b,e)(q')$, so $p'=q'$ by injectivity; components $1,2$ then
read $M(p')\bigl(p_3-q_3,\;p_4-q_4\bigr)^{\top}=0$, so $p_3=q_3$,
$p_4=q_4$.
($\Rightarrow$) Contrapositive: suppose $(b,e)(p')=(b,e)(q')$ with
$p'\ne q'$. Solve the linear system
$M(p')\,(x_3,x_4)^{\top}
=\bigl(a_{x_1}(q')-a_{x_1}(p'),\;a_{x_2}(q')-a_{x_2}(p')\bigr)^{\top}$
for $(x_3,x_4)$. Then
$\nabla f(p',x_3,x_4)=\nabla f(q',0,0)$ while
$(p',x_3,x_4)\ne(q',0,0)$: $\nabla f$ is not injective.
\end{proof}

\subsection{The trichotomy and the localization}

\begin{lemma}\label{lem:affx4}
Let $f=e_0(x_1,x_2,x_3)+x_4\,e(x_1,x_2)$ with $e\in\C[x_1,x_2]$. Then
$\det\Hess f$ is \emph{affine} in $x_4$, with
\[
[x_4]\;\det\Hess f
=-\,(e_0)_{x_3x_3}\cdot B_2(e),
\qquad
B_2(e):=e_{x_2x_2}e_{x_1}^2-2e_{x_1x_2}e_{x_1}e_{x_2}
+e_{x_1x_1}e_{x_2}^2
\]
($B_2(e)$ is the two-variable bordered Hessian of $e$).
\end{lemma}

\begin{proof}
Apply Lemma~\ref{lem:bordered} with $e_1=e$: here
$g=w:=(e_{x_1},e_{x_2},0)^{\top}$ and $K$ has only its upper-left
$2\times2$ block (equal to $\Hess_2e$) nonzero. Write $N:=P+x_4K$. Since
$w_3=0$,
\[
\det\Hess f=-w^{\top}\adj(N)\,w
=-\sum_{i,j\in\{1,2\}}w_i\,\adj(N)_{ij}\,w_j .
\]
The three relevant adjugate entries are $2\times2$ minors of $N$ using
row and column $3$:
\[
\adj(N)_{11}=N_{22}N_{33}-N_{23}^2,\quad
\adj(N)_{12}=-(N_{12}N_{33}-N_{13}N_{23}),\quad
\adj(N)_{22}=N_{11}N_{33}-N_{13}^2 .
\]
The variable $x_4$ occurs only in $N_{11},N_{12},N_{22}$ (through
$e_{x_1x_1},e_{x_1x_2},e_{x_2x_2}$), and each displayed minor contains
exactly one such entry, so each is \emph{affine} in $x_4$, with
$x_4$-coefficients $e_{x_2x_2}N_{33}$, $-e_{x_1x_2}N_{33}$,
$e_{x_1x_1}N_{33}$ respectively, where $N_{33}=(e_0)_{x_3x_3}$. Hence
$\det\Hess f$ is affine in $x_4$ and
\[
[x_4]\det\Hess f
=-(e_0)_{x_3x_3}\bigl(e_{x_2x_2}e_{x_1}^2-2e_{x_1x_2}e_{x_1}e_{x_2}
+e_{x_1x_1}e_{x_2}^2\bigr). \qedhere
\]
\end{proof}

\begin{theorem}[pivot trichotomy]\label{thm:tri}
Let $f\in\C[x_1,\dots,x_4]$ with $\det\Hess f=c\in\C^{\times}$ admit a
pivot $v$, $D_v^2f\equiv\gamma\in\C$. Then exactly one of the following
holds (the cases being indexed by the invariants $\gamma$ and
$\deg D_vf$ of the pair $(f,v)$):
\begin{enumerate}
\item[\rm(i)] $\gamma\ne0$; then $\nabla f$ is a polynomial
automorphism.
\item[\rm(ii)] $\gamma=0$ and $\deg D_vf=1$; then $\nabla f$ is a
polynomial automorphism.
\item[\rm(iii)] $\gamma=0$ and $\deg D_vf\ge2$; then, after an affine
change of coordinates,
$f=a(x_1,x_2)+x_3\,b(x_1,x_2)+x_4\,e(x_1,x_2)$ with $(b,e)$ a planar
Keller map, $\deg e=\deg D_vf$, and $\nabla f$ injective if and only if
$(b,e)$ is injective.
\end{enumerate}
\end{theorem}

\begin{proof}
Case (i) is Theorem~\ref{thm:quad}. Let $\gamma=0$ and take
pivot-adapted coordinates: $f=e_0(x')+x_4e_1(x')$ with $e_1=D_vf$. By
Lemma~\ref{lem:bordered}, $B(e_1)\equiv0$ and $e_1$ is nonconstant, so
$\deg e_1\ge1$. Case (ii), $\deg e_1=1$, is Theorem~\ref{thm:aff}.

Let $\deg e_1\ge2$. By Corollary~\ref{cor:planar} ($n=3$) there is an
affine automorphism $\chi$ of $\C^3$ with $e_1\circ\chi\in\C[x_1,x_2]$;
applying $\chi\oplus\mathrm{id}_{x_4}$ (which preserves the pivot form,
the constancy of the Hessian determinant, the degree of the pivot
coefficient, and the automorphism property) we may assume
$e_1=e(x_1,x_2)$ with $\deg e\ge2$. By Lemma~\ref{lem:affx4}, constancy
of $\det\Hess f$ forces
\[
(e_0)_{x_3x_3}\cdot B_2(e)\equiv0
\]
in the integral domain $\C[x_1,x_2,x_3]$, so one factor vanishes
identically.

We claim $B_2(e)\not\equiv0$. Otherwise, by Corollary~\ref{cor:planar}
($n=2$), after an affine change of $(x_1,x_2)$ (extended by the identity
on $(x_3,x_4)$) we may write $e=\phi(x_1)$, $\phi\in\C[s]$, with
$\deg\phi=\deg e\ge2$. Apply Lemma~\ref{lem:bordered} again with
$w=(\phi'(x_1),0,0)^{\top}$ and $K=\phi''(x_1)E_{11}$: since rows and
columns $2,3$ of $N=P+x_4K$ are those of $P$,
\[
\det\Hess f=-\phi'(x_1)^2\,\adj(N)_{11}
=-\phi'(x_1)^2\,\det\nolimits_2\Hess_{(x_2,x_3)}e_0\equiv c\in\C^{\times},
\]
and a product of two polynomials in the domain $\C[x']$ is a nonzero
constant only if both factors are nonzero constants; so $\phi'$ is a
nonzero constant and $\deg\phi=1$ --- a contradiction.

Hence $(e_0)_{x_3x_3}\equiv0$, i.e.\ $e_0=a(x_1,x_2)+x_3\,b(x_1,x_2)$,
and $f$ has the doubling form. By Theorem~\ref{thm:doubling},
$(\Jac(b,e))^2=c$ makes $(b,e)$ a planar Keller map, and
Lemma~\ref{lem:inj} gives the injectivity equivalence: case (iii). The
three cases are mutually exclusive because $\gamma$ and $\deg D_vf$ are
determined by $(f,v)$.
\end{proof}

\begin{remark}[case (ii) is not a doubling in
general]\label{rem:notdoubling}
Cases (ii) and (iii) can both be doublings (an affine $e$ with
$(e_0)_{x_3x_3}\equiv0$), but case (ii) is a genuine third class: an
earlier project draft asserted that \emph{every} affine-pivot potential
is a doubling plus an inert planar potential, and this is false. The
witness
\[
w=\tfrac12x_1^2+x_1x_2(1+x_3)+\tfrac12x_2^2(2x_3+x_3^2)+x_3x_4,
\qquad\det\Hess w=1,
\]
has the affine pivot $e_4$ with $D_{e_4}w=x_3$ of degree $1$, and its
isotropic cone $\{v\in\C^4:D_v^2w\equiv0\}$ is the single line $\C e_4$
(the coefficients of $D_v^2w=v^{\top}\Hess w\,v$ generate an ideal whose
radical contains $v_1,v_2,v_3$; Rabinowitsch certificate (W13) of
\texttt{paperG\_writeup\_checks.py}, re-verified in
\texttt{fpG\_paper\_certificate.py}). Any doubling
$a(y_1,y_2)+y_3b+y_4e$ has the $2$-plane $\mathrm{span}(e_3,e_4)$ inside
its isotropic cone, and the isotropic cone is an affine invariant (it is
carried by the linear part of the coordinate change), so $w$ is not
affinely equivalent to any doubling. Case (ii) being unconditional, this
does not affect the localization below. Case (iii) is nonempty and
nontrivial: de Bondt's anti-triangularization counterexample
$f=(x_1+x_2^2)x_3+(x_2+(x_1+x_2^2)^2)x_4$ \cite{deBondt2015} is the
doubling of an invertible planar Keller map of degree $4$. The trichotomy
is certified (Lemma~\ref{lem:affx4} fully symbolically, plus an instance
of each branch) in \texttt{pivot\_dichotomy.py}.
\end{remark}

\begin{corollary}[localization of $\JC_2$]\label{cor:headline}
The following are equivalent.
\begin{enumerate}
\item[\rm(a)] $\JC_2$ holds: every polynomial map $F:\C^2\to\C^2$ with
$\det DF\in\C^{\times}$ is a polynomial automorphism.
\item[\rm(b)] $\HC_4$ holds for every potential admitting a pivot: every
$f\in\C[x_1,\dots,x_4]$ with $\det\Hess f\in\C^{\times}$ for which some
$v\ne0$ has $D_v^2f$ constant, has $\nabla f$ a polynomial automorphism
(equivalently, polynomial Legendre transform).
\end{enumerate}
\end{corollary}

\begin{proof}
(a)$\Rightarrow$(b): let $f$ admit a pivot and apply
Theorem~\ref{thm:tri}. In cases (i) and (ii) the conclusion is
unconditional. In case (iii), $(b,e)$ is a planar Keller map, invertible
by (a), in particular injective; Lemma~\ref{lem:inj} makes $\nabla f$
injective, and Theorem~\ref{thm:inj} upgrades this to a polynomial
automorphism.

(b)$\Rightarrow$(a): let $(b,e)$ be a planar Keller map and set
$f:=x_3\,b(x_1,x_2)+x_4\,e(x_1,x_2)$. By Theorem~\ref{thm:doubling},
$\det\Hess f=(\Jac(b,e))^2\in\C^{\times}$, and $v=e_4$ is a pivot
($D_{e_4}^2f=0$). By (b), $\nabla f$ is injective; by
Lemma~\ref{lem:inj} (direction $\Rightarrow$, with $a=0$), $(b,e)$ is
injective, hence a polynomial automorphism by Theorem~\ref{thm:inj}.
\end{proof}

\begin{remark}\label{rem:sandwich}
Since $\HC_4\Rightarrow\JC_2$ (Meng's bridge), the corollary sandwiches
the pivot class exactly: $\HC_4\Rightarrow$ (b) $\Leftrightarrow\JC_2$.
So a counterexample to $\HC_4$ that is not (equivalent to) a $\JC_2$
counterexample must have no pivot in any affine coordinates, and
conversely a proof of $\HC_4$ must contain a proof of $\JC_2$ already on
the pivot class --- and, by Remark~\ref{rem:notdoubling}, must handle
case (ii) potentials that are not doublings, which
Theorem~\ref{thm:aff} does.
\end{remark}

\begin{corollary}[degree five, unconditionally]\label{cor:deg5}
Every $f\in\C[x_1,\dots,x_4]$ with $\det\Hess f\in\C^{\times}$,
$\deg f\le5$, that admits a pivot has $\nabla f$ a polynomial
automorphism.
\end{corollary}

\begin{proof}
Apply Theorem~\ref{thm:tri}. Cases (i), (ii) are unconditional. In case
(iii) the affine change preserves total degree, and
$b=\partial_{x_3}f$, $e=\partial_{x_4}f$ in the new coordinates, so
$\deg b,\deg e\le\deg f-1\le4\le100$; Moh's theorem
(Theorem~\ref{thm:moh}) makes $(b,e)$ invertible, in particular
injective; conclude by Lemma~\ref{lem:inj} and Theorem~\ref{thm:inj}.
\end{proof}

For $\deg f\le4$ the pivot hypothesis is superfluous: the companion
manuscript proves that a pivot always exists in degree $\le4$ and hence
that $\HC_4$ holds there outright \cite{companion}.

\section{What remains: potentials without a pivot}\label{sec:remaining}

Everything above is conditional on the existence of a pivot. The
remaining frontier of $\HC_4$ is therefore exactly:

\begin{problem}\label{prob:nopivot}
Does there exist $f\in\C[x_1,\dots,x_4]$ with
$\det\Hess f\in\C^{\times}$ admitting \emph{no} pivot, i.e.\ such that
for every $v\in\C^4\setminus\{0\}$ the polynomial $D_v^2f$ is
nonconstant?
\end{problem}

If the answer is negative, then by Corollary~\ref{cor:headline}
$\HC_4\Leftrightarrow\JC_2$ and the two surviving conjectures of the
family collapse into one. If the answer is positive, the no-pivot class
is the only place a counterexample to $\HC_4$ can live that would not
already refute $\JC_2$, and a search has a precisely delimited target.
Two elementary constraints on that target are worth recording.

\begin{proposition}\label{prop:pe}
Let $f\in\C[x_1,\dots,x_4]$ with $\det\Hess f\in\C^{\times}$, and write
$\Hess f=\Hess f(0)+\sum_{\alpha\ne0}x^\alpha N_\alpha$ with constant
symmetric matrices $N_\alpha$.
\begin{enumerate}
\item[\rm(a)] If $\dim\mathrm{span}\{N_\alpha\}\le3$, then $f$ admits a
pivot.
\item[\rm(b)] If $f=f_2+f_d$ with $f_2$ quadratic, $f_d$ homogeneous of
degree $d\ge3$ (one nonquadratic homogeneous part), then $f$ admits a
pivot.
\end{enumerate}
Consequently a potential without a pivot has, in every affine coordinate
system, at least four linearly independent matrices $N_\alpha$ and at
least two nonzero homogeneous parts of degree $\ge3$; by the companion's
degree-four theorem it also has degree $\ge5$.
\end{proposition}

\begin{proof}
(a) $D_v^2f=v^{\top}\Hess f\,v=v^{\top}\Hess f(0)v+\sum_\alpha x^\alpha\,
v^{\top}N_\alpha v$ is constant iff $v^{\top}N_\alpha v=0$ for all
$\alpha$, i.e.\ iff $[v]$ lies on the common zero locus in
$\mathbb{P}^3$ of $\dim\mathrm{span}\{N_\alpha\}\le3$ quadrics, which is
nonempty by the projective dimension theorem. (b) The top graded piece
of $\det\Hess f$ is $\det\Hess f_d$, of degree $4(d-2)>0$, so
$\det\Hess f_d\equiv0$; Theorem~\ref{thm:GN} ($n=4$) gives $v\ne0$ with
$D_vf_d\equiv0$, and then $D_v^2f=v^{\top}\Hess f_2\,v$ is constant.
\end{proof}

Two facts from the companion calibrate the difficulty
\cite{companion}. First, no-pivot potentials do exist one dimension up:
for every $\lambda\in\C^{\times}$ and $\mu\in\C$, the Meng--Yang
five-variable counterexample family
$\psi_{\lambda,\mu}=B+\tfrac{\lambda}{2}A^2+\mu A$ of \cite{MengYang}
admits no $v\ne0$ with $D_v^2\psi_{\lambda,\mu}$ constant (zero or not),
in any affine coordinates --- an exact Gr\"obner certificate
(\texttt{pivot\_obstruction\_d1.py}) shows the relevant cone of
candidate directions is $\{0\}$. This is the rigorous form, for the
known family, of Meng--Yang's heuristic Remark~4.2 on the
self-termination of their Schur descent, and it shows that any
localization theorem like Corollary~\ref{cor:headline} must genuinely
use dimension four. Second, the no-pivot question cannot be settled by
leading forms alone: the leading form of the Meng--Yang potential
\emph{is} a Gordan--Noether cone, and its candidate pivots are destroyed
only by lower graded pieces.

In degree five the localization has concrete bite: by
Corollary~\ref{cor:deg5}, proving that every degree-$5$ potential with
constant nonzero Hessian determinant admits a pivot would extend the
companion's degree-$4$ theorem to degree $5$ unconditionally. We have
not been able to do this.

\section{Attribution}\label{sec:attribution}

This section records provenance plainly, following the project's
adversarial prior-art audit; it overrides any stronger claim elsewhere
in the project's files.

\begin{itemize}
\item\textbf{Euler contraction (Lemma~\ref{lem:identE}): known.} It is
due to Nagaoka--Yazawa \cite{NY}: Proposition~2.3 of the arXiv version
(source label \texttt{identity1}, \S2 ``Homogeneous polynomials'')
states, for $F$ homogeneous of degree $r\ge2$ in $n$ variables,
\[
\det\bigl(-F\,H_F+s\,(\nabla F)^{\top}\nabla F\bigr)
=(-1)^{\,n-1}\,\frac{r}{r-1}\,\Bigl(s-\frac{r-1}{r}\Bigr)\,F^{\,n}\,
\det H_F ,
\]
where $\nabla F$ is a row vector and $H_F$ the Hessian. By the rank-one
expansion $\det(M+s\,uv^{\top})=\det M+s\,v^{\top}\adj(M)u$, the
$s$-coefficient of the left side is $(-1)^{n-1}F^{\,n-1}\,B(F)$;
matching with the right side and cancelling $(-1)^{n-1}F^{\,n-1}$ in the
integral domain gives Lemma~\ref{lem:identE}, and conversely the same
expansion recovers their identity from the lemma. Both directions are
machine-verified in the project certificates
\texttt{identityE\_is\_known.py} and
\texttt{gcv\_ny\_homogeneous\_converse.py}. The identity also appears in
Fox \cite{Fox2017} (in the form $(\lambda-1)U(F)=\lambda F\,H(F)$ for $F$
of positive homogeneity $\lambda$, with the same proof), and for
arbitrary $C^2$ homogeneous functions in Del Pia--Hildebrand--%
Weismantel--Zemmer \cite[Lemma 5.2]{DHWZ}, who attribute the planar case
to a 1995 manuscript of Hemmer; it is also the cone specialization of
Dolgachev's affine equation of the Hessian \cite{Dolgachev}. Cite; do
not claim.
\item\textbf{Affine equation of the Hessian (Lemma~\ref{lem:identH}):
known.} Dolgachev \cite[\S1.1.4, (1.20)--(1.21)]{Dolgachev}, with the
same proof. An earlier project draft presented it as new; that was an
attribution error, since retracted.
\item\textbf{Curvature interpretation of $B$: known.} Reilly
\cite[Prop.~4]{Reilly}; Fox \cite{Fox2017}, whose invariant $U(F)$ is
exactly $B$ and who proves $U(F)=\mathcal{K}\,|dF|^{\,n+2}$. The
underlying bordered-determinant expansion is Cauchy's.
\item\textbf{Theorem~\ref{thm:F} at $n=3$: essentially known} --- a
short corollary of de Bondt--van den Essen \cite[Thm.~3.3]{dBvdE}; see
Remark~\ref{rem:Fknown}. The Wronskian computation for the exceptional
family there is also in Fox \cite{Fox2017} (for smooth functions).
\item\textbf{Doubling (Theorem~\ref{thm:doubling}): elementary}, the
mechanism being Meng's \cite{Meng2006} (see also
\cite[Prop.~2.3]{MengYang}); only the observation that case (iii) of
the trichotomy is \emph{exactly} this shape, and the resulting $\JC_2$
equivalence, are this project's.
\item\textbf{Inputs}: Gordan--Noether \cite{GN,WdB}; Dillen
\cite{Dillen}; de Bondt \cite{deBondt2015}; Wang/Oda \cite{Wang,BCW};
Moh \cite{Moh}; Bia{\l}ynicki-Birula--Rosenlicht \cite{BBR} and
Cynk--Rusek \cite{CynkRusek}; Meng \cite{Meng2006}; the 2026
counterexamples \cite{MengYang,Gao}; Newton--Puiseux \cite{Walker} and
the extension theorems for embeddings and derivations \cite{Lang}. None
of these are ours; the Meng--Yang counterexample and Gao's determinant
claims were independently re-verified by project certificates.
\item\textbf{No prior art located} (claimed as new \emph{pending expert
review}): Theorem~\ref{thm:G} (with both proofs);
Corollary~\ref{cor:planar} under the hypothesis $B(e)\equiv0$ alone and
its sharpness (Remark~\ref{rem:dim}); Lemma~\ref{lem:affx4} and the
trichotomy Theorem~\ref{thm:tri}; the localization
Corollary~\ref{cor:headline}; Proposition~\ref{prop:pe}. Fox and Reilly
supply the framework and the identity
$B=\det\Hess e\cdot\langle\nabla e,(\Hess e)^{-1}\nabla e\rangle$, but
both work under the standing hypothesis that this quantity does
\emph{not} vanish (Fox's subject is affine spheres); the content of
Theorem~\ref{thm:G} is precisely the complementary degenerate case,
which is vacuous over $\R$ for definite Hessians, false for algebraic
functions, and has bite only for polynomials over $\C$.
\end{itemize}

\appendix

\section{Machine verification}\label{sec:verification}

All certificates are fail-closed (assert-based, nonzero exit on any
failure), in exact rational/symbolic arithmetic (sympy 1.14), in the
project directory \texttt{certificates/}. A symbolic identity checked
with fully generic coefficients in which the identity is polynomial in
those coefficients is a complete proof for all polynomials of that
shape; other machine checks are consistency checks, with the hand proofs
given above. Relevant scripts:
\texttt{fpG\_paper\_certificate.py} (this manuscript's own certificate: a
structural compile-check --- no undefined references or citations, no
duplicate labels, balanced environments and braces --- and exact
re-derivations of every algebraic identity displayed above: the
$n=1$ case, the Legendre identities $\nabla L=\psi$ and
$\Eu L-L=e\circ\psi$ for generic nondegenerate quadratics, the
structural identities $H\adj(H)g=Dg$ and $g^{\top}\adj(H)g=B$, the
flow-variant identities $\mathrm{Jac}(V)=I-H^{-1}M$ and
$\nabla(g^{\top}V)=2g-MV$, the no-positive-powers step, the Euler step
$\partial_0E|_{x_0=1}=de-x\cdot\nabla e$ and the translation invariance
of $B$ used in Theorem~\ref{thm:F}, Lemmas~\ref{lem:identE},
\ref{lem:identH}, \ref{lem:bordered}, \ref{lem:affx4}, the Schur,
doubling and collision-transfer identities, the formulas of
Theorems~\ref{thm:aff} and~\ref{thm:tri}, the witnesses of
Remarks~\ref{rem:conv}, \ref{rem:dim}, \ref{rem:strict}
and~\ref{rem:notdoubling}, and the Nagaoka--Yazawa identity with its
$s$-linear extraction); \texttt{theorem\_G.py},
\texttt{advG\_flow\_proof.py}, \texttt{pzG\_audit\_identities.py}
(structural identities of the three proofs of Theorem~\ref{thm:G};
low-degree conclusion tests; the Puiseux mechanism on an exact branch);
\texttt{theorem\_G\_n2\_deg45.py}, \texttt{pzG\_n2\_decision.py}
(exhaustive low-degree verification for $n=2$);
\texttt{tgc\_geometry\_and\_sharpness.py} (the sharpness witness of
Remark~\ref{rem:conv}); \texttt{homogenization\_identity.py}
(Lemma~\ref{lem:identH}); \texttt{cone\_leading\_form.py} and
\texttt{identityE\_is\_known.py} (Lemma~\ref{lem:identE} and its
attribution); \texttt{theorem\_F.py}, \texttt{tgc\_debondt\_chain.py},
\texttt{tgc\_beyond\_hc4.py} (Theorem~\ref{thm:F}, its second proof, and
Remark~\ref{rem:dim}); \texttt{developability.py}
(Lemma~\ref{lem:bordered}); \texttt{theoremA\_sharp.py} and
\texttt{verify\_q2\_core.py} (Theorem~\ref{thm:quad}, including fully
symbolic Schur and collision-transfer identities with non-constant $K$);
\texttt{doubling\_structure.py} (Theorem~\ref{thm:doubling} and
Lemma~\ref{lem:inj}); \texttt{pivot\_dichotomy.py}
(Lemma~\ref{lem:affx4} and an instance of each branch of
Theorem~\ref{thm:tri}); \texttt{paperG\_writeup\_checks.py} (W13: the
witness of Remark~\ref{rem:notdoubling}); \texttt{pe\_lowdeg\_pivot.py}
(Proposition~\ref{prop:pe}); and \texttt{pivot\_obstruction\_d1.py} (the
no-pivot property of the Meng--Yang family quoted in
Section~\ref{sec:remaining}).

\begin{thebibliography}{99}

\bibitem{GN} P.~Gordan, M.~Noether, \emph{Ueber die algebraischen
Formen, deren Hesse'sche Determinante identisch verschwindet}, Math.
Ann. \textbf{10} (1876), 547--568.

\bibitem{WdB} J.~Watanabe, M.~de Bondt, \emph{On the theory of
Gordan--Noether on homogeneous forms with zero Hessian},
arXiv:1703.07624 (improved version); Springer Proc. Math. Stat.
\textbf{319} (2020).

\bibitem{Dillen} F.~Dillen, \emph{Polynomials with constant Hessian
determinant}, J. Pure Appl. Algebra \textbf{71} (1991), 13--18.

\bibitem{deBondt2015} M.~de Bondt, \emph{Polynomials with constant
Hessian determinants in dimension three}, J. Pure Appl. Algebra
\textbf{219} (2015), 3743--3754; arXiv:1203.6605.

\bibitem{dBvdE} M.~de Bondt, A.~van den Essen, \emph{Singular
Hessians}, J. Algebra \textbf{282} (2004), 195--204.

\bibitem{Meng2006} G.~Meng, \emph{Legendre transform, Hessian
conjecture and tree formula}, Appl. Math. Lett. \textbf{19} (2006),
503--510; arXiv:math-ph/0308035.

\bibitem{MengYang} G.~Meng, L.~Yang, \emph{A five-variable
counterexample to the Hessian conjecture, and the low-dimensional
status of the Jacobian and Hessian conjectures}, arXiv:2607.22198.

\bibitem{Gao} S.~Gao, \emph{Counterexamples to the Jacobian conjecture
in dimensions greater than two}, arXiv:2608.00222.

\bibitem{NY} T.~Nagaoka, A.~Yazawa, \emph{Strict log-concavity of the
Kirchhoff polynomial and its applications to the strong Lefschetz
property}, J. Algebra \textbf{577} (2021), 175--202; arXiv:1904.01800,
\S2, Proposition~2.3 (source label \texttt{identity1}).

\bibitem{Fox2017} D.~J.~F. Fox, \emph{Equiaffine geometry of level sets
and ruled hypersurfaces with equiaffine mean curvature zero}, Math.
Nachr. \textbf{290} (2017), 293--320; arXiv:1503.09108.

\bibitem{Reilly} R.~C. Reilly, \emph{Affine geometry and the form of
the equation of a hypersurface}, Rocky Mountain J. Math. \textbf{16}
(1986), 553--565.

\bibitem{Dolgachev} I.~V. Dolgachev, \emph{Classical Algebraic
Geometry: A Modern View}, Cambridge Univ. Press, 2012, \S1.1.4.

\bibitem{DHWZ} A.~Del Pia, R.~Hildebrand, R.~Weismantel, K.~Zemmer,
\emph{Minimizing cubic and homogeneous polynomials over integers in the
plane}, arXiv:1408.4711.

\bibitem{Wang} S.~S.-S. Wang, \emph{A Jacobian criterion for
separability}, J. Algebra \textbf{65} (1980), 453--494.

\bibitem{BCW} H.~Bass, E.~H. Connell, D.~Wright, \emph{The Jacobian
conjecture: reduction of degree and formal expansion of the inverse},
Bull. Amer. Math. Soc. (N.S.) \textbf{7} (1982), 287--330.

\bibitem{Moh} T.~T. Moh, \emph{On the Jacobian conjecture and the
configurations of roots}, J. Reine Angew. Math. \textbf{340} (1983),
140--212.

\bibitem{BBR} A.~Bia{\l}ynicki-Birula, M.~Rosenlicht, \emph{Injective
morphisms of real algebraic varieties}, Proc. Amer. Math. Soc.
\textbf{13} (1962), 200--203.

\bibitem{CynkRusek} S.~Cynk, K.~Rusek, \emph{Injective endomorphisms of
algebraic and analytic sets}, Ann. Polon. Math. \textbf{56} (1991),
29--35.

\bibitem{Keller} O.~H. Keller, \emph{Ganze Cremona-Transformationen},
Monatsh. Math. Phys. \textbf{47} (1939), 299--306.

\bibitem{Walker} R.~J. Walker, \emph{Algebraic Curves}, Princeton Univ.
Press, 1950 (reprint Springer, 1978), Ch.~IV, \S3.

\bibitem{Lang} S.~Lang, \emph{Algebra}, revised 3rd ed., Graduate Texts
in Math. \textbf{211}, Springer, 2002, Ch.~V, \S2 and Ch.~VIII, \S5.

\bibitem{companion} HC$_4$ research project, \emph{Pivots, descents,
and the four-variable Hessian conjecture: $\HC_4$ in degree at most
four, and obstructions to the known counterexample mechanisms}, project
manuscript (\texttt{main\_result.tex}), August 2026.

\end{thebibliography}

\end{document}
